\chapter{Database-Centric Agentic System}
\label{ch:agentic}

% Sources: thesis-pack/04-architecture-agentic/agentic-core-README.md,
% agents/*.md (bus_driver.md is the reference agent), and thesis-pack/adr/.
% Stack: Ollama + a custom asyncio engine + PostgreSQL. No LangGraph, no
% PostgresSaver. Written in pure prose on purpose: no code identifiers.

This chapter describes the agentic system: how a single agent is built, how it
wakes, what it sees, what it may answer, and what the runtime does with that
answer. Chapter~\ref{ch:casestudy} already introduced \emph{who} the agents are
(Table~\ref{tab:cs:roster}) and \emph{what} they must do together
(Figure~\ref{fig:cs:storyline}). This chapter does not repeat that. It explains
\emph{how} each agent works.

The chapter describes the design as it was built and evaluated. Whether the
design helped is a question for Chapter~\ref{ch:results}.

% =====================================================================
\section{Architecture Overview}
\label{sec:ag:overview}

The system has ten agents in eight organizations. Six of them reason with a
language model: the coordinator BIS, the transport dispatcher at Hochbahn, the
two bus drivers, the police officer and the paramedic. The other four are
scripted: the disaster staff ZKD, the Augustinum and the two schools.
Section~\ref{sec:cs:agents} gives the rule behind this split. An agent that
makes no decision in the case gets no language model.

Every agent is its own small program. It runs its own loop, keeps its own state
and has its own mailbox. All ten programs are built on one shared engine, so
they wake, reason, commit and log in the same way. What differs between them
is the persona, the set of actions they may take, and a few role-specific
rules. Section~\ref{sec:ag:template} lists exactly what differs.

The agents never talk to each other directly. The only channel between them is
a PostgreSQL database, which is also the single source of truth for the whole
simulation. The database holds five kinds of record:
\begin{itemize}
  \item the current \emph{state} of every agent, one row per agent;
  \item the \emph{mailboxes}, where every message between agents is stored;
  \item the \emph{world events}, which report that a physical action has
  finished, for example that a bus has arrived;
  \item the \emph{prompt records}, which keep the exact prompt and the exact
  answer of every model call;
  \item the \emph{event log}, an append-only audit trail of everything each
  agent received and decided.
\end{itemize}
When one agent sends a message, it writes a row to the recipient's mailbox. The
database then notifies the recipient, which wakes up and reads it. The same
notification mechanism tells the observation interface that a state has
changed. Figure~\ref{fig:ag:architecture} shows this structure. Every arrow in
it passes through the database. This is what the word \emph{database-centric}
in the title of this thesis refers to. Chapter~\ref{ch:data} describes the
data layer in detail.

% Architecture of the agentic system: ten agents, PostgreSQL as the only
% channel between them. Source: agentic-core-README.md sec.13 and
% the casestudy roster. Input from inc/agentic_system.tex,
% Section "Architecture Overview".
\begin{figure}[htbp]
\centering
\begin{tikzpicture}[
  agent/.style={draw, rounded corners=2pt, fill=blue!6, minimum width=3.1cm,
    minimum height=0.62cm, align=center, font=\footnotesize},
  script/.style={draw, rounded corners=2pt, fill=gray!10, minimum width=2.8cm,
    minimum height=0.62cm, align=center, font=\footnotesize},
  group/.style={draw, dashed, rounded corners=3pt, inner sep=5pt},
  db/.style={draw, thick, rounded corners=3pt, fill=orange!8, text width=3.9cm,
    align=left, font=\footnotesize, inner sep=6pt},
  ext/.style={draw, rounded corners=2pt, fill=green!6, align=center,
    font=\footnotesize, inner sep=5pt},
  link/.style={{Stealth[length=1.8mm]}-{Stealth[length=1.8mm]}, thick, gray!70!black},
  lbl/.style={font=\scriptsize, align=center, text=gray!30!black},
]
  % --- PostgreSQL in the centre -------------------------------------
  \node[db] (db) {%
    \textbf{PostgreSQL}\\[1pt]
    \emph{single source of truth}\\[3pt]
    agent state (one row each)\\
    mailboxes\\
    world events\\
    prompt records\\
    event log (append-only)\\[3pt]
    notifications wake\\ agents and interface};

  % --- reasoning agents, left ---------------------------------------
  \node[agent, left=2.1cm of db.north west, anchor=north east, yshift=0.3cm] (bis) {BIS};
  \node[agent, below=0.18cm of bis] (hb)  {Hochbahn};
  \node[agent, below=0.18cm of hb]  (bd1) {Bus-Driver-1};
  \node[agent, below=0.18cm of bd1] (bd2) {Bus-Driver-2};
  \node[agent, below=0.18cm of bd2] (pol) {Polizei};
  \node[agent, below=0.18cm of pol] (drk) {DRK};
  \node[group, fit=(bis)(drk), label={[font=\footnotesize\bfseries]above:Reasoning agents}] (rgrp) {};

  % --- scripted agents, right ---------------------------------------
  \node[script, right=2.1cm of db.north east, anchor=north west, yshift=-0.2cm] (zkd) {ZKD};
  \node[script, below=0.18cm of zkd] (aug) {Augustinum};
  \node[script, below=0.18cm of aug] (sc1) {School SC1};
  \node[script, below=0.18cm of sc1] (sc2) {School SC2};
  \node[group, fit=(zkd)(sc2), label={[font=\footnotesize\bfseries]above:Scripted agents}] (sgrp) {};

  % --- Ollama, far left ---------------------------------------------
  \node[ext, below=0.9cm of rgrp] (llm) {Ollama\\ language model};

  % --- world / interface, below -------------------------------------
  \node[ext, below=1.6cm of db, text width=4.6cm] (ui) {%
    \textbf{World and observation interface}\\
    executes commands, reports finished actions;\\
    headless harness in the evaluation};

  % --- links (every agent link goes through the database) ------------
  \draw[link] (rgrp.east) -- node[lbl, above] {state,\\ messages,\\ log rows} (db.west |- rgrp.east);
  \draw[link] (sgrp.west) -- node[lbl, above] {messages,\\ log rows} (db.east |- sgrp.west);
  \draw[link] (llm) -- node[lbl, right] {one call\\ per wake} (rgrp.south);
  \draw[link] (ui) -- node[lbl, right, xshift=1pt] {commands out,\\ world events in} (db.south);
\end{tikzpicture}
\caption{Architecture of the agentic system. Each agent is its own program on
one shared engine. Agents never message each other directly: every message,
state change, world report and decision passes through PostgreSQL, which also
wakes the recipient. Only the six reasoning agents call the language model.}
\label{fig:ag:architecture}
\end{figure}


Vehicles are not agents. A bus, the patrol car and the ambulance are passive
resources that a person drives. The physical world is an
execution model, not a decision maker. It carries out the commands
the agents send, for example ``drive to this address'', and reports back when
each one is done. During a live run the observation interface animates this on
a map. In the evaluation a headless harness plays the same role.
% TODO(verify): "headless harness" wording, from the evaluation decision notes via the
% roster summary. Check against 10-instrument/ before the final draft.

Inference runs through Ollama. The engine around it is custom and written in
plain asynchronous Python. No agent framework is used. Section~\ref{sec:bg:agentic}
surveys such frameworks. The reason for a custom engine is that the runtime has
to own three things a general framework would not give it: time, validation and
the audit trail. Time must be owned because the agents act in a simulated clock
that keeps running while they think (Section~\ref{sec:ag:crons}). Validation
must be owned because an agent may not send a command or a message that its
role does not allow (Section~\ref{sec:ag:runtime}). The audit trail must be
owned because the evaluation reads every score from it, and the workflow
baseline has to write the very same rows. A framework that hides
these three behind its own abstractions would make the comparison in
Chapter~\ref{ch:eval} much harder to defend.

% =====================================================================
\section{The Agent Specification Template}
\label{sec:ag:template}

Every agent, in both systems, is specified along the same nine facets. The
template was fixed early in the project so that no agent is described in more
or less detail than another.

\begin{enumerate}
  \item \textbf{Context and role.} Who the agent is, where it sits, and who
  gives it instructions.
  \item \textbf{Sensing.} What the agent can perceive, and how.
  \item \textbf{Messaging.} Which message types it receives and which it may
  send.
  \item \textbf{Address book.} The concrete contacts it may reach, each with a
  reason.
  \item \textbf{Tools.} The actions it can take.
  \item \textbf{State.} What it knows and remembers, in a form that can be
  queried.
  \item \textbf{Scheduling.} How its actions are placed in time.
  \item \textbf{Instructions.} The instruction types it accepts, with example
  content.
  \item \textbf{Reasoning loop.} How it turns an instruction into a plan and a
  plan into actions. The key rule here is that there are \emph{no hard-coded
  subtasks}. The agent derives its own plan from its goal.
\end{enumerate}

The original template described scheduling as wake-ups that the agent sets for
itself. The built engine does this differently. The agent says only whether an
action should start now or as late as possible before a deadline, and the
runtime works out the exact time. Section~\ref{sec:ag:crons} explains this.

The shared engine gives every agent the whole wake cycle for free. This covers
the model call, the checking of the answer, the commit of a new step, the
replanning rules, the timing checks, the delivery of messages, the handling of
world events and every row in the event log. Each agent only has to supply six
things:
\begin{itemize}
  \item its identity, meaning the person, the address, the vehicle and the
  address book;
  \item its set of allowed actions;
  \item a way to estimate how long each action takes;
  \item its persona and its playbooks, which together form the system prompt;
  \item the way it renders its current situation into the user prompt;
  \item whether it holds a step that cannot be done in time, or lets it run
  late.
\end{itemize}
This division is deliberate. Everything that could bias the comparison between
agents, or between the two systems, lives in the shared part.

% =====================================================================
\section{The Wake Cycle}
\label{sec:ag:engine}

An agent does nothing until something wakes it. Three kinds of event wake the
language model:
\begin{itemize}
  \item a \textbf{message} arrives in the agent's mailbox;
  \item the agent's \textbf{current step finishes}, because the world has
  confirmed its last action;
  \item a \textbf{timing conflict} is found, because the step the agent just
  committed cannot be finished by its deadline.
\end{itemize}
Each of these wakes leads to exactly one model call. The answer to that call
is one decision.

\paragraph{One wake, in order.} The engine handles every wake in the same
sequence. Figure~\ref{fig:ag:engine} shows the steps, and the path beside them
that needs no model.
\begin{enumerate}
  \item It writes the incoming event to the event log.
  \item It updates the state \emph{before} the model sees it. A new message is
  put into the inbox. A finished step is moved to the history, and the next
  leg of the plan is taken from the front of the remaining plan.
  \item It takes a fingerprint of everything the model is about to see. This
  lets the evaluation tell whether two decisions were made from the same
  observation.
  \item It builds the prompt, records it, and calls the model. The prompt is
  stored before and after the call, so the interface can show that a call is
  in progress.
  \item It checks the answer against the agent's rules and counts every rule
  the answer tried to break.
  \item It commits the new step and the new plan (Section~\ref{sec:ag:replan}).
  \item It checks and sends the outgoing messages.
  \item It writes one decision row to the event log, with the reasoning, the
  step, the plan, the messages, the token counts and the latency.
\end{enumerate}

% The shared engine: how one agent handles a wake, and the model-free
% dispatch path beside it. Source: agentic-core-README.md sec.6-10 and
% agents/bus_driver.md sec.4 and sec.9. Input from inc/agentic_system.tex,
% Section "The Wake Cycle".
\begin{figure}[htbp]
\centering
\begin{tikzpicture}[
  stage/.style={draw, rounded corners=2pt, fill=blue!6, text width=4.4cm,
    minimum height=0.66cm, align=center, font=\footnotesize, inner sep=3pt},
  free/.style={draw, rounded corners=2pt, fill=gray!10, text width=3.3cm,
    minimum height=0.66cm, align=center, font=\footnotesize, inner sep=3pt},
  queue/.style={draw, thick, rounded corners=2pt, fill=yellow!12, text width=4.4cm,
    align=center, font=\footnotesize, inner sep=4pt},
  ext/.style={draw, rounded corners=2pt, fill=green!6, align=center,
    font=\footnotesize, inner sep=4pt},
  db/.style={draw, thick, rounded corners=3pt, fill=orange!8, align=center,
    font=\footnotesize, inner sep=4pt},
  dec/.style={draw, diamond, aspect=2.2, fill=gray!10, align=center,
    font=\footnotesize, inner sep=1pt},
  flow/.style={-{Stealth[length=1.8mm]}, thick},
  back/.style={-{Stealth[length=1.8mm]}, thick, dashed, red!55!black},
  link/.style={-{Stealth[length=1.6mm]}, gray!70!black},
  lbl/.style={font=\scriptsize, align=center, text=gray!30!black},
]
  % --- wake queue and the message source -----------------------------
  \node[queue] (q) {\textbf{Wake queue}\\ one wake at a time};
  \node[db, above=0.75cm of q] (mail) {New message in the mailbox};
  \draw[flow] (mail) -- node[lbl, right] {message} (q);

  % --- the model path (left column) ----------------------------------
  \node[stage, below=0.55cm of q] (p1) {1.\ Record the event in the log};
  \node[stage, below=0.28cm of p1] (p2) {2.\ Update state before the model sees it};
  \node[stage, below=0.28cm of p2] (p3) {3.\ Fingerprint the observation};
  \node[stage, below=0.28cm of p3] (p4) {4.\ Build the prompt, call the model};
  \node[stage, below=0.28cm of p4] (p5) {5.\ Check the answer against the rules};
  \node[stage, below=0.28cm of p5] (p6) {6.\ Commit step and plan\\ \emph{echo, splice or supersede}};
  \node[stage, below=0.28cm of p6] (p7) {7.\ Send the checked messages};
  \node[stage, below=0.28cm of p7] (p8) {8.\ Write one decision row};
  \foreach \a/\b in {q/p1,p1/p2,p2/p3,p3/p4,p4/p5,p5/p6,p6/p7,p7/p8}
    \draw[flow] (\a) -- (\b);

  % --- the language model --------------------------------------------
  \node[ext, left=1.0cm of p4] (llm) {Ollama\\ language model};
  \draw[link, <->] (llm) -- (p4);

  % --- the model-free path (right column) ----------------------------
  \node[free, right=1.3cm of p2] (r1) {Clock reaches an action's start time};
  \node[free, below=0.3cm of r1] (r2) {Send the action to the world\\ \emph{no model call}};
  \node[free, below=0.3cm of r2] (r3) {World carries it out and reports back};
  \node[free, below=0.3cm of r3] (r4) {Mark the action confirmed, write the real position};
  \node[dec, below=0.3cm of r4] (r5) {Last action\\ of the step?};
  \foreach \a/\b in {r1/r2,r2/r3,r3/r4,r4/r5}
    \draw[flow] (\a) -- (\b);
  \node[draw, dashed, rounded corners=3pt, inner sep=5pt, fit=(r1)(r5),
    label={[font=\footnotesize\bfseries]above:Model-free path}] (rgrp) {};
  \node[lbl, below=0.05cm of r5] {no: wait for the next report};

  % the committed step sets the action times the clock will act on
  \draw[link, dashed] (p6.east) -- node[lbl, below, pos=0.45] {action\\ times} (p6.east -| rgrp.west);

  % --- loops back into the queue -------------------------------------
  \draw[back] (r5.east) -- ++(0.55,0) |- node[lbl, above, pos=0.75] {yes: step finished} (q.east);
  \draw[back] (p6.west) -- ++(-3.9,0) |- node[lbl, above, pos=0.75] {timing conflict} (q.west);

  % --- the database ---------------------------------------------------
  \node[db, right=1.3cm of p8, yshift=0.45cm] (pg) {PostgreSQL\\ mailboxes, state, event log};
  \draw[link] (p7.east) -- (pg.west);
  \draw[link] (p8.east) -- (pg.west);
\end{tikzpicture}
\caption{The shared engine that runs every agent. A wake enters the queue and
passes through eight steps, and only step 4 calls the language model. Actions
are sent to the world on the right without the model. The model wakes again
only when a message arrives, when the last action of a step is confirmed, or
when a newly committed step cannot be finished in time.}
\label{fig:ag:engine}
\end{figure}


\paragraph{Work that needs no model.} Most of what happens in a run does not
wake the model at all. When the time for a committed action comes, the clock
sends it out to the world directly. This dispatch path makes no model call. It
only marks the action as sent and logs it. When the world reports that an
action is done, the engine marks it as confirmed and writes the vehicle's real
position back into the state. Only when the \emph{last} action of a step is
confirmed does the agent wake, with the step-finished event. Some agents also
answer certain routine messages without the model, for example a plain
acknowledgement or a position update. These \emph{fast paths} are opt-in for
each agent, because whether a message is purely informational depends on the
role that receives it. Each fast path still writes a decision row, with zero
tokens, so the event log stays complete.

\paragraph{No second chances.} The model runs at a low temperature of 0.2. If
it returns an answer that is not valid JSON, the engine does not ask again. It
records the failure and moves on. A retry would give a weak model a second
attempt that a strong model never needed, and it would hide a real difference
between models. The cost is that one bad answer can lose one decision. This is
declared in the evaluation protocol (Chapter~\ref{ch:eval}).

% =====================================================================
\section{Prompt Structure}
\label{sec:ag:prompts}

Every model call has two parts. The system prompt says who the agent is and how
it must act. The user prompt says what the situation is right now.

\paragraph{The system prompt.} It is built from three blocks, in this order.
\begin{enumerate}
  \item The \textbf{persona} says who this person is: name, age, job,
  employer, vehicle, the contacts in the address book, and the message types
  the person may send. This is the only block that is different for every
  agent.
  \item The \textbf{contract} says how the agent acts. It explains the wake
  types, the exact format of the answer, the allowed actions, how an address
  must be written, and the handover rules at the nursing home and at the
  school. The three driving agents share one contract shape. The two
  coordinators each have their own.
  \item The \textbf{playbook} is chosen by the kind of wake. A bus driver has
  three: one for planning from a new message, one for planning the next step
  after a step finishes, and one for handling a timing conflict. The
  coordinators select among more. BIS, for example, has separate playbooks for
  the first alarm, a refusal, a request for more time, a new transport plan
  and the final status of the facility.
\end{enumerate}
Each playbook is a numbered list of decision rules, followed by one worked
example. The example is set in a made-up Hamburg situation that appears in no
test scenario. This keeps the model from copying the live scenario from its own
instructions. The system prompt contains nothing specific to the scenario, so
it is the same on every wake of a run.

\paragraph{The user prompt.} It is split into fixed sections. For a bus driver
these are: the wake itself, the current location, the place where the bus is
parked, the mission brief, the mission status, the remaining plan, the current
step, the last three finished steps, and the conversations. The mission brief is
the newest dispatch order, pinned in full so it never scrolls out of view. The
conversations show the last five messages with each contact, in time order,
including the agent's own earlier messages. The agent therefore sees what it
said before, even though the model keeps no memory between calls. Every section
is always printed, with ``none'' if it is empty, so the layout never changes
from one wake to the next.

\paragraph{The answer.} The model must answer with one JSON object. It holds:
\begin{itemize}
  \item a short piece of \textbf{reasoning} in plain text, which is required;
  \item the \textbf{current step}, with a label, a deadline and an ordered
  list of actions. For each action the model says only whether it should start
  \emph{now} or \emph{by the deadline}. It never writes a clock time;
  \item the \textbf{remaining plan}, a list of future legs, each with a label
  and a deadline. If the model leaves the plan out, the old plan is kept. If it
  writes one, the new plan replaces the old one completely;
  \item the \textbf{outgoing messages}, each with a recipient, a type and a
  body;
  \item optionally, an \textbf{abort} with a reason.
\end{itemize}
The rule about the remaining plan matters for the evaluation. Because leaving
it out means ``no change'', every plan the model does write is a deliberate
revision. Appendix~\ref{app:prompts} gives the full prompts of the reference
agent.

% =====================================================================
\section{Runtime Ownership and Ground Truth}
\label{sec:ag:runtime}

The model proposes. The runtime decides what actually happens. Some state is
owned by the runtime and cannot be changed by anything the model writes. These
rules are correctness mechanisms. They are also part of the comparison, because
they limit what the agentic system can do, so they are listed here in full.

\begin{enumerate}
  \item \textbf{The runtime sets every action time.} Any time the model writes
  is thrown away and computed again from the routed duration and the fire mode.
  \item \textbf{Only allowed actions pass.} An action that is not in the
  agent's action set is dropped. The model cannot invent a new command.
  \item \textbf{Messages are checked twice.} The recipient must be in the
  address book, and the message type must be one the agent may send. A message
  that fails is dropped with a warning. It never blocks the wake, and it is
  counted as an invalid action.
  \item \textbf{Completion needs evidence.} A mission counts as complete only
  when there is no current step, no remaining plan, at least one finished step
  in the history, and no abort. An empty answer can never end a mission by
  accident.
  \item \textbf{An abort clears the plan.} After an abort, the remaining plan
  is empty, even if the same answer also wrote a new one.
  \item \textbf{History is never lost.} A step is always moved to the history
  before it is replaced. The full sequence of steps can be rebuilt from every
  run.
  \item \textbf{Each world report counts once.} A report that is delivered
  twice is recognised and ignored.
\end{enumerate}

\paragraph{Ground truth from the world.} The agent's position comes from the
world, not from the model. Each report of a finished drive carries the
vehicle's real coordinates, and the runtime writes them into the state. If a
drive is stopped half-way, for example because the agent changed its plan, the
world reports where the vehicle halted. The next prompt therefore shows where
the vehicle really is, not where the model assumed it would be. The next route
is also computed from that real position.

These rules are not an advantage for the agentic side. Several of them exist
because an earlier version failed without them. They also do not do the
agent's reasoning for it. The runtime never chooses a destination, a split of
residents or a message. It only refuses what the role does not allow, and it
keeps the record honest.

% =====================================================================
\section{Time and the Simulation Clock}
\label{sec:ag:crons}

The agents live in simulated time. Every agent runs its own clock, which moves
forward in fixed ticks. The ratio between simulated and wall-clock time is set
by a time scale. In every scored run this scale was 8.4. Only this
one value ran, so this thesis makes no claim about how the results change with
the time scale (Section~\ref{sec:eval:responsetime}).

\paragraph{Fire modes.} As Section~\ref{sec:ag:prompts} said, the model marks
each action either \emph{now} or \emph{by the deadline}. The runtime turns this
into clock times. First it estimates how long each action takes. A drive is
routed over real roads with OpenRouteService. A walk is estimated from the
straight-line distance at walking speed. Boarding and leaving a vehicle take a
fixed minute. If routing fails, the runtime falls back to a straight-line
estimate and marks the action as degraded. It never blocks the step. Then it
sets the times. A \emph{now} action starts as soon as the previous action ends.
A \emph{by the deadline} action starts as late as possible so that the whole
step still finishes on its deadline. The coordinators' actions are
administrative and take no time, so for them every action simply falls on the
deadline.

\paragraph{Timing conflicts.} After a step is committed, the runtime checks
whether any action that has not yet been sent would have had to start in the
past. A slack of 60 seconds is allowed, because the model itself takes time to
answer. If the check fails, the agent is woken with a timing conflict, together
with the size of the gap. Actions that are already sent are not checked again,
because they are facts, not plans. Changing their times on paper would not
change when they happen. Section~\ref{sec:ag:replan} describes what happens
when the conflicts repeat.

\paragraph{The clock during a model call.} A model call takes real seconds. If
the clock kept running at full speed, one slow call could move simulated time by
many minutes, and the step being planned would already be late when it arrived.
If the clock stopped completely, thinking would be free. The engine does
something in between. During a call the clock keeps running for a limited grace
window, and then freezes until the answer arrives. The window is exactly the
simulated time that the slowest allowed call can take, which is the client
timeout of 120 seconds multiplied by the time scale. Thinking therefore costs
simulated time, as it would for a real person, but the cost has an upper bound.

% =====================================================================
\section{Agent Roster}
\label{sec:ag:roster}

Table~\ref{tab:cs:roster} lists the agents. This section only describes how
the shared engine is used differently by each class of agent.

\paragraph{The bus drivers.} Both drivers use the same engine and the same
prompts. Only their identity is different. They have the richest action set:
walk, board a vehicle, leave a vehicle, drive, drive via a detour, stop, and
inform the escorts. They are the only agents that \emph{hold} a step that
cannot be done in time. A held step is not sent to the world, so the bus does
not leave while the driver negotiates. A driver who runs more than one trip
returns to the nursing home after each run, not to the depot, until the last
run is done.

\paragraph{The escorts.} The police officer and the paramedic each drive their
own vehicle as part of the convoy. The police car leads, the bus follows, and
the ambulance runs at the tail. Their contract has the same
shape as the drivers', with six physical actions. For them a timing conflict is
only a warning. The step still runs, and they tell BIS they will be late.

\paragraph{The coordinators.} BIS and Hochbahn move nothing physical. Their
actions are administrative: ask the schools for beds, request buses, order an
escort, send a dispatch, submit a report. BIS has eleven such actions and
Hochbahn seven. Each action produces a message, so BIS sends messages only
through its actions. Because nothing is dispatched to the world, their clock
does a different job. When a step's deadline arrives, it checks whether the
replies the step was waiting for have come in. If they have, the step is
finished. If not, the coordinator is woken with a timing conflict. BIS is the
only agent that decides the mission, and the only one that may call it off.
Every other agent that cannot meet a deadline asks BIS for more
time through the same request, and this includes Hochbahn.

The run schedule travels on three edges. Hochbahn sends BIS a
transport plan with every run, its passenger count and its deadline. BIS passes
the run numbers and the list of runs on to the two escorts. Hochbahn tells each
driver which run it is on and how many there are. By default bus-1 makes two
runs in sequence and bus-2 stays in reserve. Hochbahn releases the reserve for
only four reasons: a leg needs more seats than bus-1 has, bus-1 is not
available, BIS names bus-2, or the schedule fails Hochbahn's deadline test. The
baseline treats the reserve bus differently. Section~\ref{sec:val:reserve}
describes this and what it means for the comparison.

\paragraph{The scripted agents.} ZKD, the Augustinum and the two schools use
no language model. They receive a message, apply
a fixed rule and reply. Each of them still owns one fact that no other agent can
change. The Augustinum owns the count of residents who have left and the report
that the building is cleared. Each school owns the count of residents who have
arrived. ZKD owns the formal authority to close or stand down the mission.

\paragraph{The reference agent: Bus-Driver-1.} Hans Müller drives bus-1, a city
bus with 50 seats, from a depot in the north-west of Hamburg. His part of one
run goes like this.
\begin{enumerate}
  \item A dispatch order from Hochbahn arrives. Hans wakes with the
  message playbook. He accepts the mission, writes a plan of legs (drive to the
  home, load, drive to the school, unload, return), and commits the first step:
  walk to the bus, board it, drive to the Augustinum.
  \item The runtime routes the drive and sets the action times. The clock sends
  each action to the world when its time comes. No model call is needed.
  \item When the world confirms his arrival, Hans wakes with the next-step
  playbook. He tells the Augustinum he has arrived and receives the boarding
  order.
  \item After loading, he reports the bus as departed with the number of
  residents on board. His next step starts with informing the escorts. The
  runtime sends the convoy start to the police officer and the paramedic two
  minutes before the bus leaves.
  \item At the school he reports the drop-off to the school and to Hochbahn.
  If this was not the last run, his plan takes him back to the Augustinum.
\end{enumerate}
If a road closure or a late order makes a step impossible in time, Hans wakes
with the timing-conflict playbook. It gives him three choices: ask for more time
and hold, commit a step with an honest later deadline and start driving, or
abort. The playbook tells him that holding is not free, because the clock keeps
running while he waits.

% =====================================================================
\section{Replanning by Plan Rewrite}
\label{sec:ag:replan}

The system has no separate replanning module. There is no special branch that
runs ``when something goes wrong''. Instead, the model may write a new step or a
new plan on \emph{any} wake. Replanning is simply what happens when that new
step differs from the one already committed. The playbook for the wake steers
it. The timing-conflict playbook, for example, asks the model to recover from a
missed deadline. BIS has its own playbooks for a refusal and for a request for
more time.

When a new step arrives, the engine compares it with the step already
committed. There are four cases.

\begin{description}
  \item[Echo.] The model repeats the same step, with the same label, deadline
  and actions. Nothing changes. The step is not archived and not renumbered.
  \item[Splice.] The new step keeps every action that has already been sent to
  the world, and changes only the ones still waiting. The sent actions are kept
  exactly as they are. The waiting tail is replaced and timed again. For
  example, a driver who has boarded and is already driving may add a message at
  the end of the step without disturbing the drive.
  \item[Supersede.] The new step tries to change an action that has already
  been sent. That cannot be undone, so the old step is archived as aborted, and
  the new step starts from the beginning. A supersede is a change of course, not
  an edit.
  \item[Back-pressure.] Each timing conflict wakes the model again, and the
  model may answer with another late step. Without a limit, this can loop. So
  the engine stops raising timing conflicts once the reported gap no longer
  shrinks from one attempt to the next, or after three attempts in a row. The
  plan is then released to run late. A late delivery still scores, while a plan
  parked forever would not.
\end{description}

The last rule came from a real run. In one evaluation rerun, Bus-Driver-1 spent
eight timing-conflict wakes across six steps. The reported gap moved from 8 to
3, 6, 12, 7, 23, 16 and 23 minutes. The run used 869 thousand prompt tokens,
3.4 times its decision budget, and the gap never closed. Holding the step made it
worse. A held step never moves the bus, so every new step was routed again from
the same depot, and the gap kept growing.

Two points about this mechanism matter for the rest of the thesis. First, the
number of supersedes is a \emph{diagnostic}. It shows how often an agent
changed course. It is not a measure of success, and the evaluation does not
score it as one (Section~\ref{sec:eval:rubric}). Second, the ability to rewrite
the plan on any wake is the capability this architecture was built to offer.
Chapter~\ref{ch:results} reports how often the agents used it and what
followed. On the tested scenarios, runs that revised their plan more often
delivered everyone less often (Table~\ref{tab:res:commitment}). Whether the
revisions caused the losses, or followed from them, cannot be told from these
data. Section~\ref{sec:crit:mechanism} discusses this.

% =====================================================================
\section{Action Vocabulary}
\label{sec:ag:actions}

Each agent has a closed set of actions. The model may use nothing else.
\begin{itemize}
  \item \textbf{Bus drivers}, seven physical actions: walk to a place, board a
  vehicle, leave a vehicle, drive to a place, drive through waypoints, stop,
  and inform the escorts.
  \item \textbf{Police officer and paramedic}, six physical actions: the same
  set without informing the escorts, because they are the escorts.
  \item \textbf{Hochbahn}, seven administrative actions, such as dispatching a
  bus, asking BIS for more time, and reporting a run as complete.
  \item \textbf{BIS}, eleven administrative actions, such as asking the schools
  for beds, requesting buses, ordering an escort, releasing the escorts,
  submitting the mission report and asking ZKD to stand down.
  \item \textbf{Scripted agents}, none. They only reply.
\end{itemize}
Every action is checked against the set when a step is committed, and every
message against the address book and the allowed types
(Section~\ref{sec:ag:runtime}). Every rejected entry is counted. The count is
one of the evaluation's measures of decision quality
(Section~\ref{sec:eval:metrics}).

The vocabulary is also one leg of the parity argument. The workflow baseline
replaces only the two reasoning coordinators, BIS and Hochbahn
(Table~\ref{tab:val:roster}). The four field agents, meaning the two drivers,
the police officer and the paramedic, are the same programs in both systems.
Their action sets are therefore identical by construction, not by audit.
Chapter~\ref{ch:validity} sets out the full comparison, including the places
where the two systems do differ.

% =====================================================================
\section{Chapter Summary}
\label{sec:ag:summary}

The agentic system is ten small programs on one shared engine, joined only by a
PostgreSQL database. Six agents reason with a language model, and four are
scripted. A model call happens only when a message arrives, a step finishes, or
a timing conflict appears. Everything else, including sending actions to the
world, runs without the model. The model proposes a step and a plan. The
runtime owns time, checks every action and message, keeps the real position,
and writes one audit row per decision into the same event log that the baseline
uses. Replanning is not a separate module. It is any new step that differs from
the committed one, handled as an echo, a splice or a supersede, with a limit on
repeated timing conflicts. Chapter~\ref{ch:workflow} now describes the system
this design is compared against.
